% 부록: 문항을 어떻게 만들었고 실제 라벨링은 어떻게 했는지 -- 짧은 줄글 두 문단.
% 본문에서 \ref{app:prompt} 로 가리키고, 프롬프트 전문은 promptbox 로 이어서 넣는다.
% 수식·강한 주장 없이 설명조. 프리앰블에 필요한 것 없음.
\subsection{How the prompts were written and used}
\label{app:prompt}

We wrote the labelling prompts together with Claude Opus~5, revising them over
several rounds. The model was given a plain description of what the label should
mean (whether the motion after a given moment could be played back in fewer, coarser
steps without changing how the task ends), a small fixed set of demonstration
snippets from each benchmark -- camera frames at the start of a chunk, the task
instruction, a few computed facts about the chunk such as whether the gripper is
opening or closing and how fast the arm is moving, and a rough phase (reach, grasp,
carry, place, retreat) derived from the gripper signal -- and a few dozen of those
snippets that we had graded by hand. Since a VLM answers a concrete visual question
far more reliably than an abstract one, the prompt asks four to six simple checks
(``is the gripper closing on the object right now?'', ``is the object already
resting where it belongs?'') graded 1--5, some counting for and some against
compressibility, and combines them into one score. After each rewording we
re-graded the same snippets, looked at how often each check fired in each phase and
whether the hand-graded snippets came out in the right order, and kept only the
changes that clearly helped. The prompts below are the versions we stopped at.

The datasets were then labelled with Gemini~3.8~Flash through its API. Each
demonstration was cut into sixteen-frame chunks; for every chunk the model received
the camera views at the chunk's first frame, the instruction, the computed facts and
the checks, and returned one integer grade per check. The grades were combined into
the score with the weights fixed during writing, interpolated across the frames of
the chunk, and frames outside the labelled range were left out of the loss. The cost
was a few dollars per thousand chunks, and the same pipeline was used unchanged for
the real-robot datasets.
